\chapter*{Введение (задание)}
\addcontentsline{toc}{chapter}{Введение (задание)}

Цель работы — сформировать понимание методов линейного регрессионного анализа, изучить влияние штрафов $L_1$ и $L_2$ на сжатие весов признаков, а также освоить процедуру подбора гиперпараметров и оценки обобщающей способности моделей с помощью кросс-валидации.

\section*{Постановка задачи и описание набора данных}

В исследовании решается задача многомерной регрессии по прогнозированию рыночной стоимости жилой недвижимости. В качестве источника информации выбран открытый набор данных платформы PropertyFinder: \textbf{Egyptian Real Estate Listings} (доступный на Kaggle). Целевой переменной выступает стоимость объекта (\texttt{price}) в египетских фунтах (EGP).

Признаковое пространство охватывает физические и юридические параметры объектов:
\begin{itemize}
    \item \texttt{area} — общая площадь помещения в квадратных метрах;
    \item \texttt{bedrooms} — количество спальных комнат;
    \item \texttt{bathrooms} — количество оборудованных санузлов;
    \item \texttt{level} — этаж расположения квартиры;
    \item \texttt{type} — типология недвижимости (Apartment, Villa, Duplex, Penthouse, Chalet, Townhouse, Studio);
    \item \texttt{city} — территориальный район или город (New Cairo, Sheikh Zayed, 6th of October, Maadi и др.);
    \item \texttt{finishing} — уровень отделочных работ (Finished, Semi Finished, Core \& Shell, Unfinished);
    \item \texttt{furnished} — наличие мебели (Yes / No);
    \item \texttt{payment\_method} — форма расчетов (Cash, Installment, Any).
\end{itemize}

Для достижения поставленной цели решены следующие задачи:
\begin{enumerate}
    \item Загрузка данных, разведочный анализ и заполнение пропусков робастными статистиками.
    \item Кодирование категориальных параметров алгоритмом \texttt{LabelEncoder}.
    \item Разбиение выборки на обучающую и контрольную части, стандартизация числовых признаков через \texttt{StandardScaler}.
    \item Обучение базовых регрессоров: классического МНК (\texttt{LinearRegression}), модели с $L_1$-штрафом (\texttt{Lasso}) и модели с комбинированным штрафом (\texttt{ElasticNet}).
    \item Расчет метрик качества ($MSE$, $RMSE$, $MAE$, $R^2$), построение графиков сопоставления факта и прогноза, исследование разреживания весов.
    \item Поиск оптимальной силы регуляризации через \texttt{GridSearchCV} и проведение 5-кратной кросс-валидации.
\end{enumerate}

\vspace{12pt}
\noindent \textbf{Ссылка на интерактивный блокнот в Google Colab:} \\
\url{https://colab.research.google.com/github/MichaelRayven/Machine-Learning-Algorithms/blob/feat/lab4/src/lab4/lab4.ipynb}
